
\documentclass{article}
\usepackage{amsmath}
\usepackage{amssymb}
\usepackage{amsfonts}
\usepackage[parfill]{parskip}

\title{Estimating mutational variance from population-level singleton data}
\author{Sam Morris}
\date{October 2026}

\begin{document}

\maketitle

\section{Background}

We are interested in estimating $V_M$, the per-generation additive genetic variance generated by new mutations.

Our data consist of $\boldsymbol{Y}$, a vector of standardised phenotypes, and $C_i$, the number of singleton variants carried by individual $i$.

Our model is

\begin{equation}
Y_i\mid C_i \sim
\mathcal{N}\left(
\mu,\,
\sigma_e^2+C_i\sigma_{\mathrm{singleton}}^2
\right).
\end{equation}

Here, $\mu$ is the population mean phenotype, $\sigma_e^2$ is the residual variance, and $\sigma_{\mathrm{singleton}}^2$ is the variance of effects among observed singletons.

We estimate these parameters using REML.

Our aim is to estimate the variance of effects among new mutations, $\sigma_m^2$, from which mutational variance can be calculated as

\begin{equation}
V_M=U\sigma_m^2,
\end{equation}

where $U$ is the expected number of new mutations per individual per generation, for the mutation classes being analysed.

We hypothesise that
$\sigma_{\mathrm{singleton}}^2\approx\sigma_m^2$,
such that observed singletons provide a reasonable approximation to new mutations.

However, stabilising selection may preferentially remove mutations with larger effects, altering the variance of effects among observed mutations.

We therefore aim to derive an evolutionary weighting function $w(t,n;\lambda)$, where $t$ is mutation age, $n$ is allele count, and $\lambda$ describes the strength of evolutionary filtering.

We will use population-genetic theory and simulations to quantify the bias introduced by approximating $\sigma_m^2$ using $\sigma_{\mathrm{singleton}}^2$.

\section{Derivation}

\subsection{A. Population allele counts}

We aim to derive the weighting function

\begin{equation}
w(t,n)=
\frac{\mathbb{E}[\beta^2\mid t,n]}{\sigma_m^2}.
\end{equation}

We assume that the effects of new mutations are normally distributed,

\begin{equation}
\beta\sim\mathcal{N}(0,\sigma_m^2),
\end{equation}

and that stabilising selection acts on squared mutational effects according to

\begin{equation}
s=K\beta^2,
\end{equation}

where $s$ is the selection coefficient against a mutation and $K$ describes the strength of stabilising selection.

Let $x=\beta^2$. Using Bayes' rule,

\begin{equation}
p(x\mid t,n)
\propto P(n\mid t,Kx)\,p(x).
\end{equation}

To obtain the evolutionary likelihood, we use a Galton--Watson branching process.

For a rare deleterious mutation, we assume that each copy produces a Poisson-distributed number of descendant copies with mean $1-s$:

\begin{equation}
N_{g+1}\mid N_g
\sim\operatorname{Poisson}[(1-s)N_g].
\end{equation}

For this constant-population branching process, starting from $N_0=1$, the selection sensitivity of the likelihood at neutrality is

\begin{equation}
\left.
\frac{\partial}{\partial s}
\log P(N_t=n\mid s)
\right|_{s=0}
=1-n.
\end{equation}

This follows from differentiating the Poisson transition probabilities and summing across generations, where the intermediate counts cancel.

For weak selection, a first-order approximation therefore gives

\begin{equation}
P(N_t=n\mid s)
\approx
P(N_t=n\mid 0)\exp[-s(n-1)].
\end{equation}

Because $\beta$ follows a normal distribution, $x=\beta^2$ follows a Gamma distribution:

\begin{equation}
x\sim\operatorname{Gamma}
\left(\frac12,\frac{1}{2\sigma_m^2}\right),
\end{equation}

where the second parameter is the rate.

Substituting $s=Kx$ into the evolutionary likelihood and combining it with the Gamma prior gives

\begin{equation}
p(x\mid t,n)
\propto
x^{-1/2}
\exp\left[
-\left(
\frac{1}{2\sigma_m^2}+K(n-1)
\right)x
\right].
\end{equation}

The posterior is therefore another Gamma distribution, with expected value

\begin{equation}
\mathbb{E}[\beta^2\mid t,n]
\approx
\frac{\sigma_m^2}
{1+2K\sigma_m^2(n-1)}.
\end{equation}

Defining $\lambda=2K\sigma_m^2$, we obtain

\begin{equation}
w(t,n)\approx
\frac{1}{1+\lambda(n-1)}
.
\end{equation}

We note that mutation age $t$ disappears from this first-order approximation.

This result applies to the true population allele count under the specified constant-population branching process. It does not necessarily apply to allele counts observed in a finite sample. Let $n$ be the number of mutation copies in the total population and $m$ be the number of mutation copies observed in the UK Biobank sample.

We want to obtain the probability of observing $m=1$ conditional on mutation age and selection:

\begin{equation}
    P(m=1 | t, s)
\end{equation}

We can use the hypergeometric distribution to estimate the probability that we sample we sample exactly one copy of a mutation from a population.

\begin{equation}
    P(m=1 | n) = \frac{\binom{n}{1} \binom{M-n}{M_s - 1}}{\binom{M}{M_s}}
\end{equation}

The probability of observing a singleton in UKB is the sum, over all possible population allele counts, of the probability that the mutation reached that count multiplied by the probability of sampling it exactly once:


\begin{equation}
    P(m=1\mid t,s)
    =
    \sum_{n=1}^{M}
    \underbrace{P(N_t=n\mid s)}_{\text{Branching process}}
    \underbrace{P(m=1\mid n)}_{\text{Sampling}}
\end{equation}

\end{document}
